% !TEX root = ../main.tex
\section{Discussion}

Across the meta-plots the two-stage estimator reproduced the truncated fits within $6 \times 10^{-3}$ relative-frequency units and consistently delivered the lowest AICc among the three alternatives. The complete single-stage estimator remained biased near the truncation limits, whereas the truncated single-stage (1st) and two-stage complete (2sc) solutions were visually indistinguishable in Figure~\ref{fig:comparison}. These findings confirm that the two-stage workflow matches the reference truncated procedure without bespoke truncated-density implementations.

Two meta-plots (white pine--mixedwood and red pine--mixedwood) produce unusually large scale parameters in Table~\ref{tab:comparison}. Their truncated histograms are nearly flat, so the optimiser pushes the Weibull and Gamma shape parameters below one and inflates the scale (and stage-one normaliser) to hold the area to unity. Both 1st and 2sc converge on the same regime; the effect therefore reflects the data geometry rather than a weakness of the two-stage procedure.

We limited the candidate set to the two-parameter Weibull and Gamma distributions because they cover common stand structures with few parameters, both arise as special cases of the generalised gamma family, and a compact family keeps attention on the methodological swap between the truncated and two-stage estimators. A three-parameter Weibull (with location) is an alternative way to handle left truncation; we treat it as future work because it changes the parameterisation and adds flexibility that would blur the direct 1st-vs-2sc comparison in a brief format.

A companion Jupyter notebook in our GitHub repository (\url{https://github.com/UBC-FRESH/dbhdistfit-papers}) documents the plot-level tallies, weighting choices, and scripts that regenerate every figure and table, providing transparency without a separate supplement. With that reproducibility scaffold, 2sc matches or improves upon the truncated baseline for every species-group / cover-type combination, and both continue to outperform the naive 1sc approach. Because the interpretation of $b$ is application-specific, we view the two-stage workflow as a practical alternative that should be stress-tested under other truncation rules or distributional shapes before it is treated as a universal replacement. The current analysis focuses on inverse-J shaped PSP stand tables; synthetic or managed-stand shapes may reveal different strengths and should be explored in future work given the brief format.
